\section{Aplicação prática}
A interface em Streamlit e Plotly permite variar em conjunto
os coeficientes ou os polos e zeros, o ganho e o atraso.
Nyquist, Bode, Routh, lugar das raízes e resposta ao degrau
consomem as mesmas entradas e mostram a origem do modelo
racional ou exponencial usado. A exportação JSON conserva
a classificação, a convenção empregada, as margens e as amostras.
Testes de integração exercitam a troca de plantas, a atualização
dos dados e a condição de realimentação algébrica singular.

Uma bancada simulada separada usa a planta $1/(s+1)$ e um PID
com derivada filtrada, de parâmetros $K_p=2$, $K_i=1$,
$K_d=0{,}1$ e $T_f=0{,}04$ s. O modelo contínuo da função de
laço é
\begin{equation}
L_{\mathrm{PID}}(s)=
\frac{0{,}18s^2+2{,}04s+1}{0{,}04s^3+1{,}04s^2+s}e^{-s\tau}.
\end{equation}
A primeira fronteira prevista por Nyquist foi $1{,}05906$ s a
$1{,}78789$ rad/s. Em 90 s simulados a 50 Hz, os atrasos
constantes quantizados de 0,80 e 1,32 s produziram, respectivamente,
acomodação e crescimento do erro. Os RMS finais foram
$1{,}16\times10^{-9}$ e aproximadamente $3{,}76\times10^3$.
A repetição a 100 Hz preservou essa distinção. Nos cenários
com variação uniforme de até $\pm0{,}10$ s, o desvio-padrão
do atraso na fila foi próximo de 0,058 s. O comportamento
observado na simulação segue na Fig.~\ref{fig:hil}. A previsão
de Nyquist se refere ao atraso constante, enquanto a simulação
com jitter inclui discretização, fila e descarte de comandos
fora de ordem. Os resultados correspondem a um modelo linear
sem saturação. O código de comunicação para bancada permanece
sem medição em dispositivo físico.

\begin{figure*}[!t]
\centering
\figfile{0.89\textwidth}{hil-pid.pdf}
\caption{Erro no PID simulado para atrasos constantes e variáveis,
abaixo e acima da fronteira contínua prevista.}
\label{fig:hil}
\end{figure*}

\section{Discussão e limitações}
O estudo confirma que os métodos têm domínios e custos distintos.
Routh responde com uma contagem absoluta em um polinômio e usa
$O(n^2)$ operações de tabela na recorrência usual. O custo em
aritmética racional depende do crescimento de numeradores e
denominadores. Lugar das raízes resolve a dinâmica em vários
ganhos. Com matrizes densas e atribuição cúbica, o modelo de
custo típico é $O(gn^3)$ para $g$ amostras de ganho. O erro
pode crescer perto de raízes repetidas e a associação depende
da resolução da varredura.

A avaliação por Horner em $M$ pontos custa $O(Mn)$ e a
contagem de envoltórios custa $O(M)$ para uma malha dada.
O número de pontos depende do refinamento, da escala dos
polos, do atraso e da distância a $-1$. O cálculo dos polos
abertos e candidatos de cruzamento adiciona custo de raízes
polinomiais. Em Bode, a busca de margens racionais também
resolve polinômios em $\omega^2$, e no caso atrasado precisa
delimitar a banda de cruzamentos de fase. Essas diferenças
impedem atribuir superioridade por uma única medida de tempo.

Os resultados concentram-se em sistemas contínuos SISO e
em uma estação de trabalho. A validação exclui atrasos
neutros fora da cota adotada, cancelamentos numericamente
ambíguos e garantias de estabilidade sob jitter arbitrário.
Uma malha refinada pode perder uma transição entre amostras.
A estimativa de precisão para polinômios agrupados indica
que escala, tolerância e precisão de coeficientes merecem
atenção mesmo quando o estado de estabilidade concorda.
Em uma aplicação com dispositivo físico, seria necessário
medir atraso, jitter e latência de processamento, incluir
saturação e comparar séries em hardware.

\section{Conclusão}
A classificação nominal foi consistente nos quatro sistemas.
A contagem em B tornou explícito como um laço aberto
instável pode ser estabilizado pelo ganho. O zero no
semiplano direito de C apareceu no comportamento transitório,
e a comparação em D quantificou o erro das aproximações de
Padé perto da fronteira. A concordância dos métodos depende
das hipóteses do modelo, da banda e da precisão numérica
disponível. A ferramenta relaciona o diagnóstico aos
contornos, margens e raízes que o sustentam e disponibiliza
os dados para repetir a análise. Ensaios físicos e métodos
de contorno com cotas locais mais fortes são extensões do
trabalho.
